% commentaire.tex — f172r, vv. 800–837 (conversion initiale des travaux antérieurs)
\comm{801–2}{ἡμεῖς μὲν γάρ φαμεν ὑμᾶς, ὑμεῖς δ᾽ ἡμᾶς}{Molière, \italique{Impromptu de Versailles}.\citComm{\francais{Il gage que c'est moi, et moi je gage que c'est lui.}}}
\comm{804}{Ναυσιμάχης μέν ἥττων ἐστὶν Χαρμῖνος·}{Charminos, stratège athénien, qui venait de se faire battre sur mer près de Symé\nc{Cf. Thucydide, VIII, 41.}, est rapproché de Nausimaché (\italique{bataille navale}). En quoi Charminos est-il inférieur (ἥττων) à Nausimaché ? Certes par le nom -- comme Aristophane n'avait pas manqué de prévenir : après les déboires qu'il fit subir à la marine athénienne, il serait bien indigne d'en porter un de ce genre !.. mais en rester à cet unique rapprochement serait bien léger. La précision de la scholie -- Ναυσιμάχη πόρνη καὶ Σαλαβακχώ -- n'est pour une fois pas une fantaisie de Byzantin ; Nausimaché ne surpasse pas Charminos dans le \italique{combat naval} que par le nom : de par sa profession, elle devait aussi y exceller \italique{en pratique} (τἄργα) -- métaphoriquement, s'entend\nc{La langue grecque, et celle d'Aristophane en particulier, fourmille de métaphores nautiques pour rendre de façon imagée les corps-à-corps charnels ; ces images remplissent les sections 179 à 184 des \italique{Images d'Aristophane} de Taillardat et les pages 161-166 de la \italique{Maculate Muse} d'Henderson. Pour faire écho à Nausimaché, ne citons que le verbe ναυμαχεῖν usité par Aristophane aux vers 430 des \italique{Grenouilles} et 675 de \italique{Lysistrata}. Quant à τἄργα, on a déjà souligné l'acception scabreuse que peut avoir ce mot (cf. le vers 204 : ἔργα νυκτερείσια).} --, ce qui lui valut à la fois son renom et son surnom.}
\comm{805}{Κλεοφῶν χείρων πάντως δήπου Σαλαβακχοῦς.}{Salabachô est un nom générique de prostituée\nc{Hésychius (σ 88) : Σαλαβακχώ· πόρνης ὄνομα. ἀπὸ δὲ ταύτης καὶ τὰς καταφερεῖς εἰς τὰ ἀφροδίσια οὕτως ἔλεγον Ἀττικοί.}. Aristophane le mentionne aussi au vers 765 des \italique{Cavaliers} -- c'est le Paphlagonien qui parle :
\citComm{Κλέων --- τῇ μὲν δεσποίνῃ Ἀθηναίᾳ τῇ τῆς πόλεως μεδεούσῃ\\
εὔχομαι, εἰ μὲν περὶ τὸν δῆμον τὸν Ἀθηναίων γεγένημαι\\
βέλτιστος ἀνὴρ μετὰ Λυσικλέα καὶ Κύνναν καὶ Σαλαβακχώ,\\
ὥσπερ νυνὶ μηδὲν δράσας δειπνεῖν ἐν πρυτανείῳ·\\
Cléon --- \francais{À notre maîtresse Athéna, qui régit cette cité,\\
je fais cette prière ; si pour le peuple des Athéniens je me suis montré\\
le plus méritant des hommes après Lysiclès, Cynna et Salabacchô,\\
puissé-je comme maintenant, sans avoir rien fait, dîner dans le Prytanée.}}L'exégèse de ce passage nous permettra de comprendre du même coup celui qui nous intéresse.
Willems avait justement remarqué que les trois noms étaient mis ici par jeu de consonance avec les trois représentants de l'Athènes impérialiste : Périclès, Cimon et Thémistocle\nc{Platon prit les trois mêmes comme exemples d'hommes d'état honorés alors qu'ils sont en partie responsables de la crise athénienne (\italique{Gorgias}, 519a) :\citComm{ἄνευ γὰρ σωφροσύνης καὶ δικαιοσύνης λιμένων καὶ νεωρίων καὶ τειχῶν καὶ φόρων καὶ τοιούτων φλυαριῶν ἐμπεπλήκασι τὴν πόλιν.\\
sans se préoccuper de la sagesse ni de la justice, ils ont gorgé la ville de ports, d'arsenaux, de murs, de tributs \& autres niaiseries.}}. 
Pour le sens, voici ce qu'il en est : la modestie -- toujours fausse -- de Cléon, qui se place après ces trois-là, n'est que pure affectation : il fait comprendre qu'il cumule en fait les mérites des trois.
Λυσι-κλῆς, il l'est depuis sa victoire de Sphactérie : il acquit la \italique{célébrité} en \italique{déliant} la situation empêtrée du siège de Sphactérie.
Κύννα, il prétend l'être, lui qui ne cesse de proclamer être le \italique{chien} de garde de la cité\nc{Même jeu de mots avec le nom de Κύννα au vers 1032 des \italique{Guêpes} : ἀπ᾽ ὀφθαλμῶν Κύννης rappelle l'expression homérique des \italique{yeux d'un chien} (\italique{Iliade}, I, 225).}.
Σαλα-βακχώ se laisse aisément comprendre ; on rapprochera la première partie du nom d'une glose de Photius (σ 43) :\citComm{σάλα· φροντίς. οὕτως Αἰσχύλος.}On reprochait au braillard Cléon d'être un excité ; il a beau jeu de prétendre l'être en vertu de sa \italique{préoccupation}, du \italique{soin} qu'il porte à la cité. La seconde partie du nom, βακχώ, étant une allusion à l'agitation de tous les fidèles de Bacchus, et notamment les bacchantes. Suétone, dans son Περὶ βλασφημιῶν, \italique{Des termes injurieux}, cite βακχός au chapitre εἰς ἀλαζόνας, «~envers les bavards, les prétentieux~»\nc{Περὶ βλασφημιῶν, Περὶ παιδιῶν, édité en 1967 (aux Belles Lettres) par Jean Taillardat -- taillardesque édition, c'est tout dire ! βακχός est la \nieme{127} βλασφημία.}.
\textbf{Κλεοφῶν} était un démagogue dont la trop longue `carrière' politique fut brutalement interrompue en 404 av. J.-C. par une condamnation à mort. Dans les \italique{Grenouilles}, Aristophane charge ce va-t-en-guerre\nc{Vers 1504 et 1532-3.} dont la langue par trop pendue et la lippe par trop pendante n'émettaient qu'un babil barbarisant criant ses barbaresques origines\nc{Vers 678-682. L'on se reportera pour l'exégèse au toujours brillant Alphonse Willems, t. III, p. 115-117.}. (Pour un dossier complet sur Cléophon, le lecteur pourra consulter le commentaire de Fritzsche \italique{ad loc.})}
\comm{806–8}{Ἀριστομάχη, Στρατονίκη et Εὐβούλη}{sont des noms attestés au moment de la représentation des \italique{Thesmophories}\nc{Voir \italique{A Lexicon of Greek Personal Names} de P. M. Fraser et E. Matthews.}. Ces noms sont parlants, bien-sûr, mais l'allusion ; la raison pour laquelle Aristophane les cite, m'échappe. Les deux premières étaient-elles de puissantes princesses contre lesquelles les Athéniens n'auraient pas osé se frotter\nc{Στρατονίκη est porté par plusieurs membres de la famille royale macédonienne, dont l'une fut donnée en mariage à l'un des rois de Thrace Σεύθης, contemporain des \italique{Thesmophories} (voir Thucydide, II, 101 et IV, 101). Cicéron (\italique{Tusculanes}, V, 59) nous apprend que Denys lorsqu'il devint Tyran de Syracuse en 406 comptait une certaine \latin{Aristomache} parmi ses épouses.} ?}
\comm{808–9}{ἀλλ᾽ Εὐβούλης τῶν πέρυσίν τις βουλευτής ἐστιν ἀμείνων / παραδοὺς ἑτέρῳ τὴν βουλείαν; οὐδ᾽ αὐτὸς τοῦτό γε φήσεις.}{Ces deux vers ne sont qu'allusifs et ont été diversement interprétés. Ils sont d'autant plus importants que c'est peut-être le seul passage de la comédie qui puisse nous aider véritablement à dater les \italique{Thesmophories}. (Je renvoie à l'Introduction sur la question de la date de notre pièce.)
À quel événement font-ils allusion ? Deux écoles s'affrontent :\\
\ \ \ 1°) \ celle de Wilamowitz-Willems\nc{Voir la \nieme{23} note de Willems, reproduite entièrement à la fin.} qui situe la représentation des \italique{Thesmophories} aux Dionysies de 411 : le Conseil des Cinq-Cents bouleutes fut rendu tout-à-fait caduc par l'instauration d'un conseil de dix vieillards, les πρόβουλοι\nc{Voir Thucydide VIII, 1 et la \italique{Constitution d'Athènes}, XXIX ; Aristophane met en scène un Probule dans \italique{Lysistrata} (v. 387-613), jouée quelques mois auparavant -- là est la question : combien ? deux mois comme le pense Willems ?.. ou quatorze comme le pense Rogers ?}, magistrats à compétence presque illimitée qui prirent en main la direction des affaires. Willems continue : «~C'était délivrer au Conseil des Cinq-Cents un certificat d'incapacité  $[\ldots]$. Aussi, quand il sortit de charge en skirophorion (juin) 412, n'existait-il plus que de nom. [...] C'est à cette transmission annuelle de pouvoirs (παράδοσις, c'est le terme consacré), par laquelle chaque bouleute sortant cédait sa place à son successeur, que le chœur fait allusion~» ;\\
\ \ \ 2°) \ celle de Rogers -- et à laquelle je me rattache -- qui considère que la pièce est d'un an plus tardive : en mai 411, les Quatre Cents conjurés favorables à l'oligarchie se rendirent au βουλευτήριον et versèrent leur salaire aux bouleutes alors en séance. Les bouleutes se virent intimer l'ordre de se retirer, ce qu'ils firent sans même que les conjurés eussent à recourir à la violence\nc{Événement conté par Thucydide aux chapitres 69-70 de son livre VIII.}. Ainsi périt une première fois la démocratie athénienne. C'est ce manque de courage (ἀμείνων) qu'Aristophane brocarde : pas un n'aura défendu la charge qu'ils avaient reçue. Que ces piteux représentants de la démocratie fussent ceux de l'\italique{an dernier} (τῶν πέρυσίν) place donc notre pièce aux Dionysies de l'an 410.}
\comm{811–2}{οὐδ᾽ ἂν κλέψασα γυνὴ ζεύγει κατὰ πεντήκοντα τάλαντα / εἰς πόλιν ἔλθοι τῶν δημοσίων·}{Aristophane use de la même formule κατὰ πεντήκοντα τάλαντα, «~cinquante talents d'un coup~», dans les \italique{Guêpes} : Bdélycléon tance son paternel berné par les démagogues qui, non contents de dissiper l'argent de la cité, détournent le tribut des autres cités (vers 666-671) : \citComm{Φιλοκλέων --- καὶ ποῖ τρέπεται δὴ ᾽πειτα τὰ χρήματα τἄλλα;\\
Βδελυκλέων --- εἰς τούτους τοὺς “οὐχὶ προδώσω τὸν Ἀθηναίων κολοσυρτόν,\\
ἀλλὰ μαχοῦμαι περὶ τοῦ πλήθους ἀεί”. σὺ γάρ, ὦ πάτερ, αὐτοὺς\\
ἄρχειν αἱρεῖ σαυτοῦ τούτοις τοῖς ῥηματίοις περιπεφθείς.\\
κᾆθ᾽ οὗτοι μὲν δωροδοκοῦσιν \textbf{κατὰ πεντήκοντα τάλαντα}\\
ἀπὸ τῶν πόλεων ἐπαπειλοῦντες τοιαυτὶ κἀναφοβοῦντες·\\
“δώσετε τὸν φόρον, ἢ βροντήσας τὴν πόλιν ὑμῶν ἀνατρέψω.”\\
Philocléon --- \francais{Et où passe alors le reste de l'argent ?}\\
Bdélycléon --- \francais{À ceux qui nous disent: « Non, je ne trahirai point la masse athénienne, toujours je combattrai pour le populaire ! » Car c'est toi, mon père, qui te les donnes pour maîtres, enjôlé par leurs boniments. Eux cependant extorquent aux cités des cinquante talents à la fois, en les menaçant, en les terrifiant de cette sorte : « Vous paierez le tribut, ou je tonne contre votre ville et la renverse. »}}Hérodote (VI, 132-136) rapporte une histoire semblable : celle de Miltiade. Le vainqueur de Marathon, promit aux Athéniens quantité d'or s'ils lui remettaient soixante-dix vaisseaux avec troupe et argent. Les Athéniens lui fournirent le tout et Miltiade de cingler vers Paros pour en faire le siège. Il envoya un héraut réclamer 100 talents avec le même genre de menaces que sus-citées. L'expédition fut un échec et l'on intenta un procès à Miltiade pour avoir trompé la cité (τῆς Ἀθηναίων ἀπάτης εἵνεκεν). Or Miltiade se mourait d'une mauvaise blessure. 
Compte tenu des services rendus à la nation, et que c'était certainement peine perdue de transporter ce semi-cadavre jusqu'au \italique{barathre}\nc{Platon (\italique{Gorgias}, 516e) nous apprend que Miltiade devait être initialement précipité vers le \italique{barathre} (εἰς τὸ βάραθρον ἐμβαλεῖν), illustrant l'adage latin qu'il n'y a pas loin du Capitole à la roche tarpéienne. (Pour en apprendre davantage sur ce gouffre, que le lecteur se reporte au \italique{Dictionnaire des antiquités grecques et romaines}, à l'entrée \latin{Barathron}.)
Si Miltiade s'est vu gracier -- la Parque se consola néanmoins très rapidement : la grâce ne fut qu'un sursis de quelques jours --, la raison principale est peut-être qu'il bénéficiait de nombreuses relations, critère qui distingue souvent les condamnés à mort exécutés de ceux qui voient leur peine commuée, selon le mot de Gide.}, l'agonisant ne fut condamné qu'à restituer \italique{cinquante talents}.
L'histoire de Miltiade et les deux passages d'Aristophane coïncident jusque dans le chiffre du détournement !..
Miltiade, le héros de Marathon -- τὸν Μαραθῶνι, c'est son épithète, quasi-homérique, qui lui est inlassablement accolée --, profita du crédit qu'il avait auprès de Dêmos, le cajolant des souvenirs de Marathon\nc{«~Flattez-les, payez-les d'agréables mensonges : /	Quelque indignation dont leur cœur soit rempli, / Ils goberont l'appât ; vous serez leur ami.~»}, pour mener une expédition dont le but véritable était d'enflouer ses propres caisses.\\
En regard des déboires de Miltiade, la construction de l'argumentation du chœur féminin se comprend parfaitement : «~Devant Aristomaque, notre héroïne de Marathon (τὴν Μαραθῶνι), aucun d'entre vous n'entreprendrait de se battre~», et, parant d'avance à l'objection que les hommes ont aussi leur marathonien héros en la personne de Miltiade, l'on précise aussitôt : «~Une femme ayant volé cinquante talents d'un (bon) coup n'aurait pas l'outrecuidance de venir parader sur un char lors de son retour à la cité~».\\
\italique{Confer} sur ces deux vers 811-812 la note 24 d'Alphonse Willems (reproduite en annexe) : à la fête des Panathénées, la file brillante des chars était ce qui attirait le plus les regards dans la procession (Schoemann, \italique{Antiquités grecques}, t. II, p. 549).}
\comm{814}{τούτων.}{Les spectateurs.}
\comm{819–23}{}{Les trois premiers vers ne posent pas de difficulté. Et encore : à condition de garder à l'esprit les divers sens de τὰ πατρῷα, \italique{l'héritage} : désigne les \italique{biens transmis par les ancêtres} ; \italique{héritage} tout matériel que seule une femme sera à même de conserver\nc{L'homme qui avait aliéné son patrimoine, ὁ τὰ πατρῷα κατεδηδοκώς, était frappé d'ἀτιμία. Eschine, \italique{In Timarchum}, 30 ; Diogène Laërce, \italique{Solon}, I, 55. Cette loi qui n'était certainement plus observée au temps d'Eschine subsistait pour la forme, comme un vestige de l'antique règle ; il y eut toujours une δίκη κατεδηδοκέναι τὰ πατρῷα (Bekker, \italique{Anecdota}, p. 199 et 310 ; note de Fustel de Coulanges).}, mais héritage aussi immatériel, \italique{les traditions}\nc{Certains commentateurs n'ont pas manqué de noter que les Athéniennes n'auraient pas eu légalement de droit de propriété. Cette remarque aurait pu -- aurait dû -- leur mettre la puce à l'oreille : c'est bien en tant que championnes des \italique{traditions} que se posent les femmes du chœur dans ce passage.}, que les femmes perpétuent tout-de-même moins mal que les hommes (χείρους ἡμῶν εἰσιν σῴζειν). Pour le texte, voir le commentaire critique.\\
L'image ne porte pas sur l'activité filandière ; la présence simultanée de corbeilles et d'une ombrelle ne laisse pas de doute possible : Aristophane nous décrit une procession religieuse\nc{La fête des Σκίρα est évoquée quelques vers plus loin. Cf. les articles \latin{Umbella} et \latin{Skirophoria} du \italique{Dictionnaire des antiquités grecques et romaines} de Daremberg \& Saglio.}. 
Au sein de ces cortèges traditionnels, on voyait défiler les \italique{cané-}phores, qui avaient la charge des corbeilles à offrandes, et qu'elles portaient sur la tête\nc{Ces canéphores ont certainement servi de modèle aux fameuses Caryatides, portant, elles, un chapiteau sur la tête.}, quand d'autres, les \italique{skiro-}phores, les accompagnaient\nc{Voir, par exemple, le vers 1550 des \italique{Oiseaux}. Les ombrelles tenaient effectivement un rôle important dans la liturgie.} en tenant de deux mains fermes et assurées, la lourde \italique{ombrelle} consacrée, τὸ σκίρον ou τὸ σκιάδειον. 
Et pour Aristophane, sa foi, ces parades d'Athéniennes défilant cahin-caha n'étaient pas tout-à-fait indignes des traditions reçues des ancêtres (τὰ πατρῷα).}
\comm{826}{ἐκ τῶν οἴκων.}{Précisons le sens «~ἐκ τῶν οἴκων~». Ce ne sera pas du luxe. Car il ne peut être question de rompre des lances \italique{à l'extérieur des maisons} plutôt qu'\italique{à l'intérieur} ? Devrai-je aussi commenter les commentaires qui voient en cette panoplie militaire le fameux héritage paternel, τὰ πατρῷα, que les fils prodigues indignes dissiperaient -- ô scandale !
Au garde-à-vous, je félicite ces brillants universitaires pour leurs éclaircissements : sans eux, je crois que les douceurs martiales de la poésie d'Aristophane m'auraient échappé. Notre poète roulant des variations sur le thème : «~Tu seras un soldat, mon fils !~», c'est à vous tirer des larmes (de guerre, bien-sûr !)\\
Οἶκος, ce mot rare et qui renvoie à des réalités n'existant pas ailleurs qu'en Grèce classique a tout-de-même donné le terme οἰκο-νομία, qui me semble être la bonne gestion d'une \italique{maison}, perçue comme un \italique{fonds} ou un \italique{capital}. «~ἐκ τῶν οἴκων~» désigne tout ce que l'on \italique{tire de} ce fonds, ou de ce capital. La bonne équation est évidemment celle-ci : τὰ πατρῷα $ = $ ἐκ τῶν οἴκων.
La même tournure avec le même sens se trouve chez Xénophon dans son ... \italique{Économique} (XX, 21) :\citComm{τὰ οὖν συντρίβοντα τοὺς οἴκους πολὺ μᾶλλον ταῦτά ἐστιν ἢ αἱ λίαν ἀνεπιστημοσύναι. τὸ γὰρ τὰς μὲν δαπάνας χωρεῖν ἐντελεῖς \textbf{ἐκ τῶν οἴκων}, τὰ δὲ ἔργα μὴ τελεῖσθαι λυσιτελούντως πρὸς τὴν δαπάνην\\
Voilà qui ruine une maison beaucoup plus sûrement qu'une excessive ignorance : d'un côté, dépenser complètement \italique{ce qu'on tire des maisons}, de l'autre des travaux qui ne sont pas profitables à la dépense}Traduisez donc :\citComm{τοῖς δ᾽ ἡμετέροις ἀνδράσι τούτοις\\
ἀπόλωλεν μὲν πολλοῖς ὁ κανὼν\\
ἐκ τῶν οἴκων αὐτῇ λόγχῃ\nc{Écho au proverbe : αὐτῷ κανῷ... sauf que le \italique{canon} est remplacé par la \italique{lance}. Autre sexe, autre proverbe :\citComm{ἦν γάρ τις αἶνος, ὡς γυναιξὶ μὲν τέχναι\\
μέλουσι, \textbf{λόγχῃ δ᾽ ἄνδρες} εὐστοχώτεροι.\\
εἰ γὰρ δόλοισιν ἦν τὸ νικητήριον,\\
ἡμεῖς ἂν ἀνδρῶν εἴχομεν τυραννίδα.\\
Il existe en effet un proverbe qui dit qu'à la femme les artifices reviennent, tandis que \italique{les hommes sont plus habiles à atteindre le but par la lance}. Si le prix de la victoire allait à la ruse, c'est nous qui régnerions sur les hommes.}(Euripide, fr. 6 \italique{Danaé} Jouan-van Looy $=$ fr. 321 Kannicht $=$ fr. 413 Mette.)}.\\
Du fait de nos ci-devant maris (elles les désignent),\\
la lance a anéanti pour beaucoup\\
ce qu'on tirait de nos maisons -- lance comprise.}}
\comm{827–9}{πολλοῖς δ᾽ ἑτέροις ἀπὸ τῶν ὤμων / ἐν ταῖς στρατιαῖς / ἔρριπται τὸ σκιάδειον.}{Un bouclier est une \italique{ombrelle}. Cette métaphore unique est très claire : un bouclier est rond (cf. κύκλος, Taillardat § 642) et bombé comme une ombrelle. Mais on peut se demander si l'argot des soldats athéniens ne connaissait aussi pas cette métaphore unique\nc{À propos de θύλακος (§268), on [l'«~on~», ou l'«~om~» en question n'est autre que le grand Taillardat !] a déjà relevé une métaphore ironique qui a dû appartenir à l'argot des soldats athéniens. Cette tendance à ironiser sur ses armes a toujours existé : on sait que les soldats de César appelaient, par plaisanterie, \latin{cippi} et \latin{lilia} certaines défenses qu'ils construisaient devant Alésia (\italique{Bell. Gall.} 7, 73, 4 et 8). C'est de la même façon que dans l'argot militaire français, un fusil peut s'appeler une \italique{arbalète}, une \italique{seringue}, une \italique{canne à pêche}, etc. (voir A. Dauzat, \italique{L'argot de la guerre}\expo{2} (1919).)} : le bouclier protège des coups comme l'ombrelle du soleil, cf. Σ : τὴν ἀσπίδα λέγει· σκιάζει\nc{Σκιάζειν, \italique{couvrir} : Hésiod. \italique{Théog.} 716 ; Hdt. 6, 117, etc.} γὰρ τὸν ἄνδρα. Cette hypothèse est malheureusement invérifiable.
(Note de Taillardat, §643. L'utilisation deux vers plus haut de κανών comme synonyme de \italique{lance}, λόγχη, pourrait relever de ce même argot militaire dont Taillardat fait la très astucieuse hypothèse.)\\
La charge soldatesque des hommes ? elle leur passe par-dessus la tête et les épaules lorsqu'ils en délestent leur bouclier (ἀπὸ τῶν ὤμων ἔρριπται τὸ σκιάδειον), bien encombrant il est vrai lorsqu'il faut prendre la poudre d'escampette en plein milieu des combats... 
Mica, la femme de Cléonyme -- Cléonyme qui semble avoir fait école avec sa célèbre ῥίψασπις (voir le commentaire au vers 605) -- devait probablement à ce moment être pointée, au moins des yeux.}
\comm{834}{Στηνίοισι καὶ Σκίροις.}{Fêtes athéniennes réservées aux femmes. C'est notamment aux Sciries que fut élaboré le projet politique de Praxagora (\italique{Assemblée des femmes}, vers 18).}
\comm{833–7}{ταξίαρχον ἢ στρατηγόν ... ἢ τριήραρχον πονηρὸν ἢ κυβερνήτην κακόν}{Les exemples d'Aristophane sont tirés de l'actualité la plus brûlante. Les mauvais triérarques et commandants de vaisseaux sont ceux responsables de la défaite navale d'Érétrie à l'été 411, défaite qui plongea les Athéniens dans une consternation, nous rapporte Thucydide (VIII, 95 et 96), plus grande encore que celle consécutive au désastre de Sicile. Défaite cinglante qui ne rattrape pas aux yeux d'Aristophane la récente et lointaine victoire de Κυνὸς Σῆμα (VIII, 104-108) dans l'Hellespont. (Cette victoire devait être bien amère à Aristophane, puisque son contemporain Eupolis y mourut. J'y reviendrai.)}
